\section{特殊案件的管辖}
\begin{enumerate}
  \item \yi 先识别凭什么管。保护、属人、拟制领土、普遍（条约管辖）。\er 找连接点。
  \item \kd{外国人犯罪}。
        \begin{enumerate}
          \item 普遍管辖。\zd{抓获}地、\zd{登陆}地、\zd{入境地}。
          \item 保护管辖。保护中国人的管辖方法。\zd{被告登陆和两境，被害可由现前管}（离境前、现居住地）。
        \end{enumerate}
  \item \kd{外国人在国外针对中国人犯罪}。保护管辖。\zd{被告登陆和两境；被害可由现前管}。
  \item \kd{中国人在国外犯罪}。属人管辖。
        \begin{enumerate}
          \item 使领馆内。司法豁免回原籍（含主管单位所在地）。
          \item 使领馆外。属人管辖。\zd{被告登陆、入境地，双方都有（现}居住地，离境\zd{前}居住地）\zd{管}。\al{陈世峰}。
        \end{enumerate}
  \item \kd{交通工具}。不看逃窜、在哪被抓。
        \begin{enumerate}
          \item 飞机。领域外 \yjt 飞机\zd{最初降落地}。
          \item 中国船舶。内水领海 \yjt 普通地域 + 登陆。领域外 \yjt \zd{初泊、登陆、入境地}。
          \item 列车。
                \begin{enumerate}
                  \item 国际列车。\zd{没有协议，始发停靠}。
                  \item 国内列车。\yi 运行中被抓 \yjt \zd{前方停靠始终管}。\er 非运行中 \yjt 乘警地（负责该列车乘务的铁路公安机关）+ 车站抓获（车站所在地）。
                \end{enumerate}
        \end{enumerate}
  \item \kd{服刑/执行期间}。\zd{漏罪首选原审地，新罪首选服刑地，异地服刑犯罪地，脱逃抓到且发现}。
\end{enumerate}
